%  LaTeX support: latex@mdpi.com 
%  For support, please attach all files needed for compiling as well as the log file, and specify your operating system, LaTeX version, and LaTeX editor.

%=================================================================
\documentclass[knowledge,article,submit,pdftex,moreauthors]{Definitions/mdpi} 

%=================================================================

% MDPI internal commands
\firstpage{1} 
\makeatletter 
\setcounter{page}{\@firstpage} 
\makeatother
\pubvolume{1}
\issuenum{1}
\articlenumber{0}
\pubyear{2023}
\copyrightyear{2023}
%\externaleditor{Academic Editor: Firstname Lastname}
\datereceived{12 February 2023} 
%\daterevised{} % Only for the journal Acoustics
\dateaccepted{} 
\datepublished{} 
\hreflink{https://doi.org/} % If needed use \linebreak
%\doinum{}

%=================================================================
% Add packages and commands here. The following packages are loaded in our class file: fontenc, inputenc, calc, indentfirst, fancyhdr, graphicx, epstopdf, lastpage, ifthen, lineno, float, amsmath, setspace, enumitem, mathpazo, booktabs, titlesec, etoolbox, tabto, xcolor, soul, multirow, microtype, tikz, totcount, changepage, attrib, upgreek, cleveref, amsthm, hyphenat, natbib, hyperref, footmisc, url, geometry, newfloat, caption

\graphicspath{{figures/}} % path to folder with figures
\usepackage{subcaption}
\usepackage{listings}
\usepackage{xcolor}
\usepackage{gensymb} % for \textdegree
\usepackage{tabularx}

%=================================================================
% Full title of the paper (Capitalized)
\Title{Desertification in the Khartoum Region of Sudan Analysed from the Satellite-Derived Vegetation Indices}

% MDPI internal command: Title for citation in the left column
\TitleCitation{Desertification in the Khartoum Region of Sudan Analysed from the Satellite-Derived Vegetation Indices}

% Author Orchid ID: enter ID or remove command
\newcommand{\orcidauthorA}{0000-0002-5759-1089} % Polina Add \orcidA{} behind the author's name
\newcommand{\orcidauthorB}{0000-0002-6461-1551} % Olivier Add \orcidB{} behind the author's name

% Authors, for the paper (add full first names)
\Author{Polina Lemenkova $^{1}$*\orcidA{} and Olivier Debeir $^{1}$\orcidB{}}

%\longauthorlist{yes}

% MDPI internal command: Authors, for metadata in PDF
\AuthorNames{Polina Lemenkova and Olivier Debeir}

% MDPI internal command: Authors, for citation in the left column
\AuthorCitation{Lemenkova, P.; Debeir, O.}
% If this is a Chicago style journal: Lastname, Firstname, Firstname Lastname, and Firstname Lastname.

% Affiliations / Addresses (Add [1] after \address if there is only one affiliation.)
\address{%
$^{1}$ \quad Laboratory of Image Synthesis and Analysis (LISA), École polytechnique de Bruxelles (Brussels Faculty of Engineering), Université Libre de Bruxelles (ULB). Building L, Campus du Solbosch, ULB -- LISA CP165/57, Avenue Franklin D. Roosevelt 50, Brussels 1050, Belgium.
}

% Contact information of the corresponding author
\corres{Correspondence: polina.lemenkova@ulb.be; Tel.: +32 471 86 04 59 (P.L.)}

% Current address and/or shared authorship
%\firstnote{Current address: Affiliation 3.} 

\abstract{Desertification in Sudan-Sahel region of Africa has contributed to the growing interest in remote sensing data analysis. The vegetation indices (VIs) illustrate the important parameters of crop health and plant greenness. They present a crucial indicator of vegetation coverage, which reflect ecosystem structure and functional features of canopy that is essential for environmental and climate analytics. The Earth observing satellite sensors providing remote sensing data to support research enable to improve our understanding of links between the environmental parameters and vegetation responses to hydrological and climate processes leading to desertification. In this study, we calculated several VIs to compare the dynamics of the VIs in the region of confluence of Blue Nile and White Nile in Khartoum, southern Sudan. The data were compared for the periods of 2013, 2018 and 2022. We tested five following VIs for the Landsat 8-9 OLI/TIRS images covering Khartoum area: Normalized Difference Vegetation Index (NDVI), Normalized Difference Water Index (NDWI), Red-Edge Chlorophyll Vegetation Index (RECl), Modified Soil-Adjusted Vegetation Index (MSAVI) and Optimized Soil Adjusted Vegetation Index (OSAVI). The demonstrated dynamics in the changes of the VIs collected with 5-4 year interval may thus be considered as indicator to reveal climate–vegetation relationships. Using programming methods for computing and visualization of vegetation indices presents new challenges for image processing, such as the Landsat 8-9 OLI/TIRS scenes. Therefore, in this paper, we demonstrated the technical use of R scripts and packages 'raster' and 'terra'  for spatial data processing to automate image analysis and mapping. We demonstrated that the climatic drivers of desertification affecting plant growth and health varied for images taken on different years along with changed precipitation.}

% Keywords
\keyword{Remote sensing; Sudan; Khartoum; Nile River; image analysis; image processing; information retrieval; programming; Africa; computer vision} 

\PACS{91.10.Da; 91.10.Jf; 91.10.Sp; 91.10.Xa; 96.25.Vt; 91.10.Fc; 95.40.+s; 95.75.Qr; 95.75.Rs; 42.68.Wt}
\MSC{86A30; 86-08; 86A99; 86A04}
\JEL{Y91; Q20; Q24; Q23; Q3; Q01; R11; O44; O13; Q5; Q51; Q55; N57; C6; C61}

%%%%%%%%%%%%%%%%%%%%%%%%%%%%%%%%%%%%%%%%%%
\begin{document}

\section{Introduction}

Desertification is one of the most destructive climate-related issues in Sudan, especially for its southern and central regions. Desertification results from various factors including climate change and human activities. The drivers of climatic variations include annual decrease of rainfalls, soil salinization and related deficit in moisture \cite{DAWELBAIT201245}. Located southwards of Sahara Desert in the Sudan-Sahel region of Africa, between the latitudes 8° and 23°N (see Figure \ref{fig01}), Sudan experiences the especially great risk of desertification. Within the area of the country, the most endangered region prone to desertification lies between 12° to 18°N, including the area surrounding the capital of Sudan -- Khartoum. Climate factors contributing to the hazard of desertification in Sudan include the seasonality of the rainfall patterns \cite{MOHAMED20223265}, the instability and irregularity of wet and dry seasons \cite{HULME198689}, the variability of precipitation \cite{Elagib,Ayoub} and the overall water deficit \cite{akhtar1994methods}. This results in high evapotranspiration: 1,900 mm in the extreme north near Sahara and over 2,500 mm in southern Sahel and savannah \cite{ElKarouri1986}. Moreover, flat topography of the country exaggerates the desertification problems of Sudan and increases the effects from dust and sand storms often in the northern regions of the country \cite{khiry2006mapping,9635812,7948717,7381366}. Specifically, this includes the effects from the sandstorms and creeping sands and dunes from the areas of Sahara, Nubian and Libyan deserts, Figure \ref{fig01}.

\begin{figure}[H]
\centering
	\includegraphics[width=14.0 cm]{Topo_SD.jpg}
	\caption{Topographic map of Sudan. Mapping software: Generic Mapping Tools (GMT) scripting toolset. Data source: GEBCO/SRTM. Cartography source: authors.
	\label{fig01}}
\end{figure}

Social factors of desertification include a complex chain of processes related to land management, agriculture policymaking and governance \cite{KOBAYASHI2020257}. Although agriculture and livestock raising play an important role in the economy of Sudan as the major sources of livelihood for inhabitants \cite{sulieman2010expansion}, unsustainable agricultural expansion and farming activities also contribute to the increase of desertification \cite{Ibrahim}. For instance, uncontrolled agriculture results in soil erosion, compaction and nutrient depletion, changed irrigation patters and increase in pesticides \cite{Nesser} contribute to land cover changes. Moreover, overgrazing of cattle and pastoral activities \cite{Sulieman} modify land cover types and landscape structure \cite{Ahmed}. 

Furthermore, disorganized use of grazing lands and the sedentarization of the nomadic tribes also mentioned by \cite{Akhtar} as factors contributing to desertification. Despite existing efforts on water resources conservation and rational exploitation of resources \cite{MOHAMEDALI2003237}, the environmental sustainability in Sudan region remains a challenge. In turn, social consequences of desertification in Sudan result in the increased possibility of the severe drought and associated famine due to affected food production and soil degradation in the Sahelian zone of the Sudan \cite{Subramaniam}. A long-term reconstruction and modelling revealed the environmental changes already presented in the past, as well as gave reference to the links between the human occupation and their effects on changes in river discharge and occasional floods, which are in turn forced by climatic changes \cite{HONEGGER2015141}.

Recent clime change studies on Sudan confirm rising temperatures, declining rainfall, decreased relative humidity, as well as changes in patterns of cloudiness, atmospheric radiation and evaporation \cite{Alvi}. These processes trigger environmental degradation and intensification of desertification, especially in vulnerable zones of central and southern Sahelian zone of Sudan \cite{Pearce,YANG2022106213} and savannah \cite{Falkenmark}. As a consequence, this results in deforestation, soil desiccation, changed or affected land cover types \cite{Biro} and decreased biodiversity. The latter includes, for instance, the extinction of wild species in areas of oases where they dominated. A further consequence of land cover changes may lead further to the imbalance of ecosystems with spread of invasive species and disturbance of natural habitat patterns \cite{Fuller}. Recent studies reported \cite{SALIH2017S31} frequent drought periods during the last three decades. The severe droughts of the 1970s and 1980s, caused by the decrease in rainfall, resulted in several cycles of famine in Sudan with affected farmers and population. Regional studies report cases of drought and desertification in diverse regions of Sudan, e.g., chronic drought during the past 15 years in Sahel region of Darfur \cite{FULLER198739}, central Sudan \cite{Salih2022}, or eastern Sudan \cite{TARIKU2021144863}. The morphodynamic processes indicating desertification can be used for operative monitoring and detecting the vulnerable regions. Apart from the vegetation patterns, these may include soil erosion, alterations of relief specifically in the sandy regions of western Sudan \cite{Bakhit}. Finally, the cumulative effect of all these factors is mentioned by \cite{Darkoh} who noted the overlap of the effects from the unpredictable and severe droughts, desiccation or aridification and dryland degradation or desertification in the dryland regions of Africa. 

Further, desertification leads to many negative environmental consequences such as land degradation in arid and semi-arid regions of Sudan \cite{Khiry,ELNASHAR2022152925}. For instance, deep ploughing of the surface increased the susceptibility of soil to wind erosion, which has led to a severe decline in fertility and contributed to the formation of sand dunes \cite{Abdi}. More examples of the effects of the desertification include reduced biological productivity \cite{KHALIFA2018790,Jahelnabi} and decrease in plant biomass, increase of the extremely dry desert areas. 

\begin{figure}[H]
\makebox[1.0\linewidth][r]{%
	\begin{subfigure}[b]{.6\textwidth}
		\centering
			\includegraphics[width=0.90\textwidth]{Google_Earth_1.jpg}
		\caption{}
	\end{subfigure}%
	\begin{subfigure}[b]{.6\textwidth}
		\centering
		\includegraphics[width=0.90\textwidth]{Google_Earth_2.jpg}
		\caption{}
	\end{subfigure}%
}
\caption{The confluence of the White and Blue Nile in Khartoum, visualized based on the map and aerial images from the Google Earth: (\textbf{a}) Google Earth Map, (\textbf{b}) Google Earth Image.}\label{fig02}
\end{figure}

The specifics of the Khartoum region is deeply connected to the Nile River system and the confluence of the White Nile and Blue Nile, Figure \ref{fig02}. The Nile River is a central water artery of the country of which 70\% belongs to its basin \cite{Melesse}. The hydrological processes of the Nile River result in diverse impacts on the environment, including occasional floods, soil erosion and succession of vegetation on the islands and banks \cite{Halwagy} and accumulation of sedimentation \cite{BILLI201012}, growth of aquatic plants and weeds \cite{PACINI2016242}. The Nile also provides a specific places for wetland habitat and biodiversity of rare species. The immense contribution of the Nile Basin for social development of Sudan can be illustrated by its crucial role for irrigated agricultural and hydropower. Thus, the Nile serves as a central source of life in Sudan providing water resources for about 85\% of Sudanese population and irrigated agriculture for mixed cultivation of beans, sorghum, and sugarcane \cite{Kauetal2023}. Additionally, it is a transport waterway connecting the country with neighbouring territories. Five from the total six Cataracts of Nile are located within Sudan (Figure \ref{fig01}), while only one (the First Cataract) is situated in Egypt near the Aswan Dam. 

Further, the importance of the Nile River Basin can be illustrated by the fact that roughly 2/3 of its territory lies within Sudan and its major tributaries -- the While Nile and Blue Nile -- confluence near its capital, the Khartoum \cite{Hamad}. Intense practice of urban agriculture in Khartoum aims at meeting the increasing food demand of the rapidly growing population \cite{SCHUMACHER2009399} which necessarily results in additional pressure on the environment on the one hand, and the gradual increase of the built-up area on the other. The confluence of the Blue Nile and White Nile forming its major artery of the Greater Nile near the Khartoum city (Figure \ref{fig02}) results in the high concentration of population in this district which, as a consequence, leads to the increased environmental pressure in the region. The other important tributaries are the Atbara River (sometimes referred as Red Nile) with high mining activities related to gold and copper mining and cement production \cite{SASMAZ2020106539}. 

General climatic trends are necessarily reflected in the hydrology of the Nile river which demonstrates high levels of interannual and interdecadal variability in river flow records \cite{CONWAY200599}. For instance, the increasing trend in natural water storage variation of the Nile during 2013–2016 is detected over the northern region compared to the southern Sahel \cite{ABDELMALIK2019103596}. Such fluctuations in the hydrology of the main water artery necessarily affect the vegetation coverage in the banks and surrounding areas close to Nile. Besides the direct effects from climate, the patterns of the vegetation of Khartoum Province are deeply connected to the regional geologic setting \cite{ELSHEIKH201182}, the geomorpic structure and dominating soil types that control the distribution of the specific plants \cite{Mahmoud}.

In this paper, we introduce the application of R libraries for remote sensing data processing, and, in discussing their technical details, focus on the calculating several satellite-derived vegetation indices based on Landsat 8-9 OLI/TIRS images collected on Khartoum city around the confluence of the White Nile and the Blue Nile. The method can be applied to other types of imagery besides the Landsat (e.g., MODIS, AVHRR, Sentinel). The goal in the presented series of the analysed images on Khartoum is to model the changes in vegetation patterns and greenness which can be analysed based on the comparison of the time series and to visualise the dynamics in vegetation coverage. The standard GIS-based approaches for this problem are based on using the traditional software, for instance, ERDAS Imagine \cite{AHMED201829}, and involve many manual operations and routine work \cite{Hawash,YOUSSEF2020103777,6621869,Aldoma}. Instead of using such conventional approaches,  we combined several libraries of R such as 'raster' and 'terra' with auxiliary packages in a principled way to obtain a better workflow for image processing and analysis.

%%%%%%%%%%%%%%%%%%%%%%%%%%%%%%%%%%%%%%%%%%

\newpage
\section{Materials and Methods}

The importance of the use of the remote sensing data and advanced methods of image processing is raised in earlier studies \cite{ELBASTAWESY2017252,SOSNOWSKI201651,GORSEVSKI201364}. For examples, Landsat images are known for reliable source of data for relatively accurate technique for classifying time-series and detecting forest and land cover types \cite{WILSON2002385,Lemenkova202227,4293065}, computing vegetation indices \cite{5416935,Lemenkova202229} and specifically desertification \cite{RIVERAMARIN2022104829} to show the advance or retreat of arid areas using analysis of the satellite images.


%----------- subsection ----------->
\subsection{Data}

\begin{figure}[H]
\makebox[0.90\linewidth][r]{%
	\begin{subfigure}[b]{.40\textwidth}
		\centering
			\includegraphics[width=0.90\textwidth]{Fig_02a.jpg}
		\caption{2013}
	\end{subfigure}%
	\begin{subfigure}[b]{.40\textwidth}
		\centering
			\includegraphics[width=0.90\textwidth]{Fig_02b.jpg}
		\caption{2018}
	\end{subfigure}%
	\begin{subfigure}[b]{.40\textwidth}
		\centering
			\includegraphics[width=0.90\textwidth]{Fig_02c.jpg}
		\caption{2022}
	\end{subfigure}%
}
\caption{Landsat 8 OLI/TIRS image of White and Blue Nile in Khartoum, Sudan, in natural RGB colours in December: (\textbf{a}) 20.12.2013, (\textbf{b}) 18.12.2018, (\textbf{c}) 21.12.2022}\label{fig:03}
\end{figure}

%----------- subsection ----------->
\subsection{Methodology}

We processed the satellite images using libraries of R programming language \cite{RCoreTeam} using adapted methods from the spatial packages. The map in Figure \ref{fig01} has been plotted using Generic Mapping Tools (GMT) scripting toolset Version 6.1.1. \cite{Wessel} by applying the existing cartographic workflow \cite{Lemenkova202301,Lemenkova202230,Lemenkova202228}.

%----------- subsection ----------->
\subsection{Subsection}

%%%%%%%%%%%%%%%%%%%%%%%%%%%%%%%%%%%%%%%%%%
\section{Results}

%%%%%%%%%%%%%%%%%%%%%%%%%%%%%%%%%%%%%%%%%%
\section{Discussion}

%%%%%%%%%%%%%%%%%%%%%%%%%%%%%%%%%%%%%%%%%%
\section{Conclusions}

%%%%%%%%%%%%%%%%%%%%%%%%%%%%%%%%%%%%%%%%%%
\authorcontributions{Supervision, conceptualization, methodology, software, resources, funding acquisition, and project administration, O.D.; writing---original draft preparation, methodology, software, data curation, visualization, formal analysis, validation, writing---review and editing, and investigation, P.L. All authors have read and agreed to the published version of the manuscript.}

\funding{The publication was funded by the Editorial Office of Knowledge, Multidisciplinary Digital Publishing Institute (MDPI), by providing 100\% discount for the APC of this manuscript. This project was supported by the Federal Public Planning Service Science Policy or Belgian Science Policy Office, Federal Science Policy - BELSPO (B2/202/P2/SEISMOSTORM).}

\institutionalreview{Not applicable.}

\informedconsent{Not applicable.}

\dataavailability{Not applicable} 

\acknowledgments{The authors thank the anonymous reviewers of Land for reading, suggestions and comments that improved an earlier version of this manuscript.}

\conflictsofinterest{The authors declare no conflict of interest.} 

\abbreviations{Abbreviations}{
The following abbreviations are used in this manuscript:\\

\noindent 
\begin{tabular}{@{}ll}
MDPI & Multidisciplinary Digital Publishing Institute\\
DOAJ & Directory of open access journals\\
LD & Linear dichroism
\end{tabular}
}

%%%%%%%%%%%%%%%%%%%%%%%%%%%%%%%%%%%%%%%%%%
%% Optional
\appendixtitles{no} % Leave argument "no" if all appendix headings stay EMPTY (then no dot is printed after "Appendix A"). If the appendix sections contain a heading then change the argument to "yes".
\appendixstart
\appendix
\section[\appendixname~\thesection]{}
\subsection[\appendixname~\thesubsection]{}
The appendix 

%%%%%%%%%%%%%%%%%%%%%%%%%%%%%%%%%%%%%%%%%%
\begin{adjustwidth}{-\extralength}{0cm}
%\printendnotes[custom] % Un-comment to print a list of endnotes

\reftitle{References}

%=====================================
% References, variant A: external bibliography
%=====================================
\bibliography{Sudan.bib}
%\nocite{*}

%%%%%%%%%%%%%%%%%%%%%%%%%%%%%%%%%%%%%%%%%%
\end{adjustwidth}
\end{document}
